% Copyright Ben Paul Wise. All Rights Reserved.
% ----------------------------------------------
% Part II -- Developer manual / API. First-drafted part (it fixes the
% terminology). Sources: the headers under include/ (each engine's contract
% comment is the authoritative statement), CLAUDE.md, CMakeLists.txt.
% ----------------------------------------------
\chapter{Developer Manual}
\label{ch:manual}

This chapter is written for an engineer who wants to \emph{use} the library
(pose a problem, pick an engine, interpret the result) or \emph{extend} it
(add an engine, an NCP function, or a structured linear solver). It fixes the
vocabulary the rest of the report uses. All code resides in the namespace
\code{VINCP}; all mathematical statements are deferred to
Chapter~\ref{ch:math}, and appear here only as far as an API contract
needs them.

\section{Getting started}
\label{sec:getting-started}

The library is C++20 over Eigen, and builds on both Windows~11 (MSVC) and
Debian Linux (GCC/Clang); nothing platform-restricted is used. Eigen is
header-only: there is nothing to compile or install for it. The build locates
it through the CMake cache variable \code{EIGEN3\_INCLUDE\_DIR}, which must
point at the directory that \emph{contains} the \code{Eigen/} folder (the
checkout root, not the \code{Eigen} subfolder); on Debian,
\code{/usr/include/eigen3} after \code{apt install libeigen3-dev}.

\begin{lstlisting}
cmake -S . -B build -DCMAKE_BUILD_TYPE=Release
cmake --build build
ctest --test-dir build --output-on-failure
\end{lstlisting}

Tests are GoogleTest suites, fetched at configure time by CMake
\code{FetchContent} at a fixed release tag and compiled by the same
toolchain with the same flags (on MSVC, \code{gtest\_force\_shared\_crt} is
forced on to match the dynamic CRT). A single suite runs by ctest name
(\code{ctest -R lcp\_psd\_test}) or directly as an executable. Three
aggregate custom targets run groups during the \emph{build} step:
\code{run\_all\_tests}, \code{run\_engine\_tests}, and
\code{run\_network\_tests}; in an IDE they are \emph{built}, not run (a
custom target has no executable of its own).

The directory layout, in dependency order (arrows point one way, no cycles):

\begin{center}
\begin{tabular}{ll}
\toprule
\code{include/}, \code{lib/} & the core library \code{vincp}: types, engines, drivers \\
\code{network/}              & the flow-planning layer \code{vincpnet} (Section~\ref{sec:network}) \\
\code{test/}                 & GoogleTest suites for the core library \\
\code{src/}                  & runnable demos (printout, not pass/fail checks) \\
\code{doc/}                  & papers (local copies), design notes, this report \\
\bottomrule
\end{tabular}
\end{center}

The library target is built with \code{/W4} (MSVC) or \code{-Wall -Wextra}
elsewhere, and the working convention is that the build stays free of
warnings; Eigen's headers are included as \code{SYSTEM} and exempt.

\section{Core types}
\label{sec:core-types}

One header, \code{vincp.hpp}, defines the domain vocabulary. It is included
by everything and includes nothing of the library itself. It is also the only
header that names the \code{Eigen} namespace: it pulls \code{VectorXd},
\code{MatrixXd}, \code{Index}, and the factorization types into
\code{VINCP}, so the Eigen dependency stays localized and reversible.

\subsection{The problem: \code{VIModel}}

The library's target problem is the \emph{mixed complementarity problem}
(equivalently, a variational inequality over \K): find $z = (x, y)$ with the
free block $x \in \Rn$ and the nonnegative block $y \in \Rm$ such that
\[
  H(x, y) = 0,
  \qquad
  0 \le G(x, y) \perpc y \ge 0 .
\]

\begin{lstlisting}[language=C++]
struct VIModel {
  Index n = 0;   // dimension of the free block x
  Index m = 0;   // dimension of the non-negative block y
  function<VectorXd(const VectorXd& x, const VectorXd& y)> H;  // -> R^n
  function<VectorXd(const VectorXd& x, const VectorXd& y)> G;  // -> R^m
};
\end{lstlisting}

The free function \code{evaluateF(model, z)} splits $z$ into
$x = \code{z.head(n)}$, $y = \code{z.tail(m)}$, stacks $(H, G)$, validates
block lengths, and throws on a non-finite value. When the map is more natural
as a single field $F:\mathbb{R}^{n+m} \to \mathbb{R}^{n+m}$,
\code{makeVIModel(n, m, F)} builds the split form from it. Affine problems
(the engines of Section~\ref{sec:engines} that take \code{(M, q)} directly)
do not need a \code{VIModel} at all.

\subsection{The result: \code{VIResult}}

Every solver in the library --- leaf engine, driver, or composition ---
returns the same result type:

\begin{lstlisting}[language=C++]
struct VIResult {
  VectorXd z;             // solution vector
  double   residual  = 0.0;   // SQUARED residual norm at termination
  int      iter      = 0;     // iterations (outer, for a composite solver)
  bool     converged = false; // residual fell below tolerance within the cap
  int      innerIters = 0;    // total inner iterations of a composite; 0 for a leaf
};
\end{lstlisting}

Two contract points deserve emphasis because they are easy to get wrong:

\begin{itemize}
\item \code{residual} is the squared Euclidean norm of the \emph{natural
  residual}: the standard defect vector of complementarity theory, which is
  zero exactly at solutions. On our $K$ it reads componentwise, $H(x,y)$ on
  the free rows and $\min(y_i, G_i)$ on the nonnegative rows, so it measures
  equation violation, infeasibility, and non-complementarity in one vector
  (the formal projection form and its properties are
  Section~\ref{sec:a-problem}). Every tolerance in the library
  (\code{magTol}, \code{outerTol}, \ldots) is likewise a \emph{squared}
  norm, a deliberate, library-wide convention
  (Section~\ref{sec:contracts}).
\item \code{converged = false} is a defined outcome, not an error. A solver
  that stalls --- on a failed line search, a vanishing merit gradient, a
  collapsed interior-point step length, or the Josephy--Newton no-progress
  cutoff --- returns its iterate with \code{converged = false} rather than
  throwing. Two solvers additionally guarantee \emph{which} iterate:
  \code{semismoothNewtonSolve} and \code{alternatingChainSolve} return the
  best point visited (in the natural-residual sense), not the last one,
  because their nonmonotone strategies may deliberately move above the best
  before stalling.
\end{itemize}

What \code{iter} counts differs by engine (an iteration of \code{bsHe94b} is
one back-substitution; of \code{mehrotraIpm}, one LU factorization; of
\code{solveVI}, one linearization with a full inner solve); the per-engine
meanings are collected in Section~\ref{sec:engines} and revisited when
comparing engines in Chapter~\ref{ch:testing}.

\subsection{The feasible set: \code{Projector}}

The set $K$ enters the projection engines only through a
\code{Projector = function<VectorXd(const VectorXd\&)>}: given $v$, return
the Euclidean projection of $v$ onto $K$. Ready-made:

\begin{itemize}
\item \code{projectNonnegative} --- onto $\mathbb{R}^{d}_{+}$
  (componentwise $\max(v, 0)$);
\item \code{makeMixedProjector(numFree)} --- onto
  $\mathbb{R}^{\code{numFree}} \times \mathbb{R}^{d-\code{numFree}}_{+}$,
  the set matching the \code{VIModel} split;
\item \code{makeEllipsoidProjector(radii, tol, iterMax)}
  (\code{ellipsoidprojector.hpp}) --- onto an axis-aligned solid ellipsoid,
  by a bracketed root-find on the KKT multiplier. It exists both for its own
  use and as a working demonstration that the projection engines are generic
  in $K$.
\end{itemize}

Not every engine accepts an arbitrary $K$: the interior-point method and the
semismooth Newton solver compile the mixed structure
$\K$ into their algebra and take \code{numFree} (or the \code{VIModel})
instead of a \code{Projector}. Each engine's subsection states which
convention it uses.

\subsection{Logging hooks}

All iterative engines share
\code{IterationLogger = function<void(int iter, int iterMax, double mag,
double magTol)>}, called every \code{iterFreq} iterations (\code{iterFreq}
$\le 0$ disables logging; an empty logger is never called). The
Josephy--Newton driver's \code{OuterLogger} and the alternating chain's
\code{ChainStageLogger} follow the same opt-in pattern at their own
granularities. Loggers are the intended instrumentation mechanism:
production runs stay silent, and a long diagnostic run attaches a flushed
per-iteration heartbeat (a lesson recorded in Chapter~\ref{ch:testing}:
under a piped test runner, unflushed output from a killed run is lost).

\section{Library-wide contracts}
\label{sec:contracts}

Three conventions hold across every component; code that joins the library
is expected to preserve them.

\paragraph{Squared norms.} Every residual and every tolerance is a
\emph{squared} Euclidean norm ($\code{dot(e,e)}$, never $\code{norm(e)}$).
The convention matches the Octave reference implementations the library
descends from and the tolerances validated against them; because it is
uniform, cross-engine comparisons (and this report's tables) are directly
comparable. Do not convert a squared residual to a plain norm: halving the
exponent of every published tolerance would silently change every
acceptance test.

\paragraph{Bad values throw; stalls return.} The library never silently
substitutes a result. The dividing line:
\begin{itemize}
\item A \emph{bad value} --- a NaN residual, divergence past the guard
  ($\code{mag} \ge 1 + \code{divergenceFactor} \cdot \code{initialMag}$), a
  non-finite linear solve, an inconsistent Jacobian shape --- throws
  \code{std::runtime\_error}. Caller mistakes (dimension mismatches,
  parameters outside their ranges, unset functors) throw
  \code{std::invalid\_argument}, validated up front.
\item A \emph{stall} --- no descent direction, a failed line search, an
  interior-point step length below \code{stallStep}, the Josephy--Newton
  no-progress cutoff --- is a legitimate outcome at a real iterate: the
  solver returns it with \code{converged = false}. The distinction has
  practical consequences: engine \emph{compositions}
  (Section~\ref{sec:altchain}) treat a stalled stage as ``switch strategy''
  rather than ``abort,'' and a throw where a stall belongs forces every
  wrapper to catch and second-guess it.
\end{itemize}

\paragraph{One result type, one convention per extension point.} Every
solver returns \code{VIResult}; every swappable component is a
\code{std::function} extension point with a documented contract
(\code{Projector}, \code{InnerSolver}, \code{NewtonSolverFactory},
\code{NcpFunctionPair}, \code{JacobianFn}, \code{StageSolver}). An
extension point is preferred over a boolean flag wherever a third
implementation is plausible.

\section{Engine catalog}
\label{sec:engines}

The library deliberately offers a catalog of engines rather than one
universal algorithm; Chapter~\ref{ch:testing} documents, on real problems,
why no single entry suffices. Engines divide into three groups by the
problem form they accept:

\begin{itemize}
\item \textbf{Affine engines} solve the linear VI / mixed LCP (linear
  complementarity problem, $F(z) = Mz + q$) given \code{(M, q)} directly:
  \code{dHan06}, \code{bsHe94b}, \code{solodovSvaiter},
  \code{chainedSolodovHe} (over any \code{Projector}), and
  \code{mehrotraIpm} (mixed/orthant structure only).
\item \textbf{Nonlinear engines} take a \code{VIModel}:
  \code{solveVI} (Josephy--Newton over any \code{Projector}) and
  \code{semismoothNewtonSolve} (mixed structure compiled in).
\item \textbf{Compositions} assemble engines:
  \code{chainedSolodovHe} at the affine level,
  \code{alternatingChainSolve} at the model level.
\end{itemize}

The affine engines share one calling convention ---
\code{engine(x0, M, q, Pr, magTol, iterMax, iterFreq, params, logger)} ---
so they are interchangeable wherever that signature is expected; only the
params struct differs. Each subsection below gives the engine's problem
class, its distinguishing mechanics, its parameters, its failure modes, and
when to choose it.

\subsection{\code{dHan06} --- self-adaptive projection}
\label{sec:dhan06}

Han's 2006 self-adaptive projection method for the linear VI over any
projector-defined $K$. Convergence theory requires \emph{monotone} $F$:
for affine problems, $M$ positive semidefinite, the mathematically
well-behaved class to which the optimality systems of convex programs
belong (definitions in Section~\ref{sec:a-problem}). Each iteration factors
$(I + \beta_k M)$ with a self-adaptive $\beta_k$; \code{DHan06Params}
exposes the relaxation factor \code{gamma} $\in (0,2)$ (default 1.6), Han's
constant \code{mu} (1.05), the initial \code{beta0}, the $\tau$-schedule
(\code{tau0}, \code{tauN}), the divergence guard, and the diagnostic switch
\code{adaptBeta}.

Two practical notes. First, \code{beta0 = 1.0} (the default, a deliberate
departure from the Octave source's $0.5$) makes the \emph{first} step's
metric $(I + \beta_0 M) = (M + I)$, which is exactly \code{bsHe94b}'s
proven-contractive metric, so the method starts stably even when a
linearization is indefinite and adapts from there. Second, because $\beta_k$
changes, the factorization is redone every iteration: on equal iteration
counts \code{dHan06} has been observed an order of magnitude slower in wall
time than \code{bsHe94b}. With \code{adaptBeta = false} the method is
identical to \code{bsHe94b}, which isolates the self-adaptive rule for
A/B comparison.

\emph{Choose it} when the self-adaptive metric justifies its
refactorization cost; in practice this is rare, and the engine serves
mainly as the reference implementation of its paper and as the comparison
baseline.

\subsection{\code{bsHe94b} --- fixed-metric projection-contraction}
\label{sec:bshe94b}

He's 1994 projection-contraction method (the paper's equation~16 search
direction) for the linear VI over any $K$. It factors $(M + I)$ once
and reuses the factorization every iteration, so the per-iteration cost is
one back-substitution plus projections. \code{BsHe94bParams} is minimal:
\code{gamma} $\in (0,2)$ (default 1.6) and the divergence guard.

This is the library's standard engine for well-conditioned monotone affine
problems, and the finishing stage of \code{chainedSolodovHe}. Its known
weakness is geometric: on near-degenerate problems (wide, flat optimal
faces from near-tied costs) the contraction rate approaches 1 and the
iteration count grows by orders of magnitude --- the failure that motivated
the interior-point engine (Chapters~\ref{ch:testing} and \ref{ch:math}).

\emph{Choose it} for monotone affine problems of moderate conditioning,
for \emph{warm-started} re-solves (runs started from a previous, nearby
solution, which projection engines exploit fully), and wherever
an arbitrary projector-defined $K$ (e.g.\ an ellipsoid) rules out the
structural engines.

\subsection{\code{solodovSvaiter} --- hyperplane double-projection}
\label{sec:ss}

The Solodov--Svaiter (1999) projection method: one projection produces the
residual $r = x - P_K(x - \mu F(x))$; an Armijo-style search along $r$ (a
backtracking search that shrinks the trial step until a sufficient-decrease
test passes --- the pattern recurs throughout the library) locates a
hyperplane separating the iterate from the solution set; the next iterate
is the projection onto that hyperplane, re-projected onto $K$.
\code{SolodovSvaiterParams}: \code{mu} (1.0), \code{sigma} (0.5),
\code{gamma} (0.5), \code{maxBacktracks} (60), divergence guard.

It has two distinctions: it is matrix-free ($M$ enters only through
products; nothing is factored), and it converges globally for merely
pseudomonotone continuous $F$, the weakest hypothesis of any engine here.
Its equally important limitation: whenever $\|F(x^\ast)\| > 0$,
the hyperplane step collapses like $\|r\|^2/\|F\|^2$ and the squared
residual decays only as $O(1/\sqrt{k})$. It is a globalizer rather than a
finisher --- effective for moving from an arbitrary start into a solution's
neighborhood, too slow to refine an answer to tight tolerance --- and its
own tests assert only globalization-grade accuracy by design.

\emph{Choose it} as the safe first phase from a bad start on a problem whose
monotonicity is in doubt, or where memory forbids factorization, and pair
it with a finisher.

\subsection{\code{chainedSolodovHe} --- globalize, then contract}
\label{sec:chained}

The two-phase affine composition: \code{solodovSvaiter} to a loose target
(\code{roughMagTol}, default $10^{-4}$ squared, capped at
\code{roughIterMax}), then \code{bsHe94b} warm-started from its iterate and
run to the caller's \code{magTol}. Phase 1 reaching its iteration cap is
not an error; producing a warm start is its purpose.
\code{ChainedSolverParams} nests the two phases' params. Accounting follows
the composite convention: \code{iter} is phase 2's count, \code{innerIters}
the total across both. This chain is the small-scale ancestor of the
model-level alternating chain (Section~\ref{sec:altchain}): each method
plays only its strong regime.

\subsection{\code{fbsHyz04} --- forward-backward splitting}
\label{sec:fbs}

The He--Yuan--Zhang (2004) forward-backward splitting scheme
\cite{hyz2004}, ported from the author's earlier \code{pmedemo}
implementation (whose adaptive step rule it preserves). Unlike every other
engine at this level it takes the \emph{nonlinear} field directly ---
\code{fbsHyz04(x0, F, Pr, magTol, iterMax, iterFreq, params, logger)} with
\code{F} a \code{VectorField} --- so it solves the general monotone VI with
no Josephy--Newton wrapper and no Jacobian of any kind: two $F$ evaluations
and two projections per iteration, nothing factored, any projector-defined
$K$. The step $\gamma$ adapts online from a smoothed secant estimate of the
Lipschitz constant (\code{FbsHyz04Params}: \code{gamma0}, \code{theta},
\code{lcFactor}, defaults from \code{pmedemo}). Convergence theory needs
monotone Lipschitz $F$. An affine adapter, \code{makeFbsHyz04Solver}, fits
the \code{InnerSolver} extension point.

Two documented safeguards go beyond the 2004 recipe: a conservative
sampled-secant preflight for $\gamma_0$ (automatic when \code{gamma0}
$\le 0$), and Tseng-style step backtracking, which extends robustness to
merely locally-Lipschitz fields; on well-behaved fields neither has any
effect.

\emph{Choose it} for nonlinear monotone VIs where Jacobians are unavailable
or unaffordable (black-box $F$, memory limits) or where a simple robust
iterator matters more than a fast tail. Its historical role: it solved the
strategic allocation game in the PME reference implementation, and its
library port independently confirms that game's reference equilibrium
(Chapter~\ref{ch:testing}).

\subsection{\code{mehrotraIpm} --- predictor-corrector interior point}
\label{sec:ipm}

An infeasible primal-dual predictor-corrector interior-point method
(Mehrotra-style path following) for the \emph{monotone mixed LCP} over
$\K$: with $z = (x, y)$ and $M, q$ partitioned conformally,
$(Mz+q)_x = 0$ and $0 \le y \perpc s = (Mz+q)_y \ge 0$. \code{numFree = 0}
gives the pure LCP. This is exactly the KKT-of-convex-QP class: the
first-order optimality system (Karush--Kuhn--Tucker conditions) of a convex
quadratic program, with $x$ the primal and free variables and $y$ the
inequality multipliers.

\begin{lstlisting}[language=C++]
VIResult mehrotraIpm(const MatrixXd& M, const VectorXd& q, Index numFree,
                     double magTol, int iterMax, int iterFreq,
                     const MehrotraIpmParams& params = {},
                     const IterationLogger& logger = {},
                     const NewtonSolverFactory& newtonFactory = {});
\end{lstlisting}

Structural departures from the affine convention, each deliberate:
\begin{itemize}
\item \textbf{No start vector.} An interior-point method cannot use a caller
  iterate (it needs $y, s$ strictly positive near the central path); it
  starts from a data-scaled interior point. Consequently it cannot
  warm-start, which matters for re-solve loops.
\item \textbf{No \code{Projector}.} $K$ is structurally
  $\mathbb{R}^{\code{numFree}} \times \mathbb{R}^{m}_{+}$; pairing the
  adapter (Section~\ref{sec:jn}) with any other $K$ silently solves the
  wrong problem, so its header forbids it.
\item \textbf{One factorization per iteration}, shared by the predictor and
  corrector solves; \code{iterMax} therefore counts factorizations, and
  sensible caps are in the low hundreds, not the tens of thousands of the
  projection engines.
\end{itemize}

The engine exists because its iteration count ($\sim$10--40) is essentially
independent of dimension \emph{and of degeneracy} of the solution set:
the central path ends at the analytic center of the optimal face, so the
near-tied instances that stall projection-contraction cost it only a few
extra iterations (mechanism in Chapter~\ref{ch:math}; evidence in
Chapter~\ref{ch:testing}).

\code{MehrotraIpmParams}: \code{tauFraction} (0.995 fraction-to-boundary),
centering clamps \code{sigmaMin}/\code{sigmaMax}, \code{stallStep}
($10^{-10}$; a damped step below this returns with
\code{converged = false} --- a stall, not an error), \code{regEpsilon}
($10^{-8}$, the sticky free-block diagonal rescue applied only if a Newton
solve returns non-finite values), the divergence guard, and
\code{newtonCheckTol} (below).

\paragraph{The \code{NewtonSolverFactory} extension point.} The
per-iteration linear algebra is swappable, layered as in OOQP: the
interior-point recipe is fixed, the Newton solve is not.

\begin{lstlisting}[language=C++]
using NewtonSolve = function<VectorXd(const VectorXd& rhs)>;
using NewtonSolverFactory =
    function<NewtonSolve(const VectorXd& sOverY, double freeRegularization)>;
\end{lstlisting}

The factory is called once per iteration and must return a solver for
$K = M + \mathrm{blockdiag}(\code{freeRegularization} \cdot I_n,\,
\mathrm{diag}(\code{sOverY}))$, where $M$ and $n$ are the data the
\emph{caller} bound when constructing the factory. An empty factory means
the built-in dense-LU implementation (bit-for-bit the engine's historical
behavior). Because the engine cannot see the factory's $K$,
\code{newtonCheckTol > 0} enables a development-mode drift guard: one
$O(d^2)$ matvec verifies $\|Kd - \mathit{rhs}\|^2$ after every solve and
throws on disagreement. \code{freeRegularization} carries the singularity
rescue through the interface. The flow-planning layer's
\code{makeFlowNewtonFactory} (Section~\ref{sec:network}) is the production
example: it replaces the dense $O(d^3)$ factorization with per-sink
Sherman--Morrison and a small Schur complement, the change that turned a
30-minute solve into 4.5 seconds.

\emph{Choose it} for monotone affine problems that are degenerate,
near-degenerate, or large --- especially when problem structure admits a
custom Newton factory --- and accept that it cannot warm-start and that its
$K$ is fixed by structure.

\subsection{\code{solveVI} --- the Josephy--Newton driver}
\label{sec:jn}

The outer driver for the \emph{nonlinear} VI over a projector-defined $K$:
each step linearizes $F$ at $z_k$ (4th-order finite-difference Jacobian,
\code{fdjacobian.hpp}), hands the affine VI $(M = J(z_k),\
q = F(z_k) - J(z_k) z_k)$ over the \emph{same} $K$ to an inner solver, and
damps the returned step with an Armijo line search on the natural-map merit
(without damping, the undamped step exhibits period-2 chatter at non-smooth
solutions). The projector carries the mixed structure, so there is
deliberately no Schur complement and no elimination of the free block.

\begin{lstlisting}[language=C++]
using InnerSolver = function<VIResult(const VectorXd& x0, const MatrixXd& M,
                                      const VectorXd& q, const Projector& Pr)>;
VIResult solveVI(const VIModel& model, const VectorXd& z0,
                 const InnerSolver& innerSolver,
                 const JosephyNewtonParams& params = {},
                 const OuterLogger& logger = {},
                 const Projector& projector = {});   // empty => mixed projector
\end{lstlisting}

The inner solver is a parameter, not a hard-wired choice. Adapters bind each
concrete engine's controls into the \code{InnerSolver} shape:
\code{makeDHan06Solver}, \code{makeBsHe94bSolver},
\code{makeSolodovSvaiterSolver}, \code{makeChainedSolver} (all generic in
$K$), and \code{makeMehrotraIpmSolver(numFree, ...)} --- which ignores both
the start (no warm start) and the projector (structural $K$), and must only
be used where the mixed/orthant $K$ is the set actually intended. Inner
tolerances are set in the adapter, outer controls in
\code{JosephyNewtonParams}: \code{outerTol}, \code{outerIterMax}, the
Jacobian step \code{fdStepRel}, the Armijo block, a diagnostic
\code{logInnerDefiniteness} (prints the smallest eigenvalue of the
symmetrized linearization --- the monotonicity probe), and the no-progress
cutoff \code{stallIterMax}/\code{stallRelDecrease}: after
\code{stallIterMax} \emph{consecutive} outer iterations without relative
improvement of the best residual, the driver stops and reports
\code{converged = false} \emph{before} spending another Jacobian and inner
solve. The cutoff is off by default.

Two tuning facts with evidence behind them: the outer residual floor scales
linearly with the \emph{inner} tolerance (tighten \code{innerMagTol}, not
the Jacobian, to reach a tighter \code{outerTol}); and the FD Jacobian order
does not bind that floor (2nd- and 4th-order give identical iterates).

\paragraph{The forcing sequence.} Binding one tight inner tolerance up
front wastes the early iterations: far from the solution, the linearization
is about to move anyway, and solving it to $10^{-12}$ yields no benefit. The
\emph{forcing-sequence} overload of \code{solveVI} takes an
\code{InnerSolverFactory} --- \code{function<InnerSolver(double
innerMagTol)>}, a one-line lambda over any \code{make*Solver} adapter ---
and rebuilds the inner solver each outer iteration at
\[
  \mathit{innerTol}_k \;=\;
  \mathrm{clamp}\bigl(\code{forcingRatio}\cdot \mathit{residual}_k,\;
  \code{forcingFloor},\; \code{forcingCap}\bigr),
\]
loose while the residual is large, tight near the solution (defaults
$10^{-2}$ / $10^{-14}$ / $5\times 10^{-4}$, matching the reference Octave
scripts' \code{min(5e-4, n0/100)} rule; theory in
Chapter~\ref{ch:math}). The observable effect is the reference profile:
thousands of inner iterations early, around a hundred late, with no loss of
final accuracy.

\paragraph{The one-call entry.} For callers who want a single call with no
decisions:
\begin{lstlisting}[language=C++]
VIResult r = solveVIVanilla(model, z0);   // optional: outerTol, outerIterMax
\end{lstlisting}
binds the \code{bsHe94b} inner under the forcing sequence with every other
control at its default. One deliberate caveat, stated in its header: it has
no basin control --- on problems with multiple equilibria it converges
quickly to whichever equilibrium its trajectory enters; deliberate
selection among equilibria is the alternating chain's task
(Section~\ref{sec:altchain}).

\emph{Choose it} for smooth nonlinear VIs over any projector-defined $K$,
with the inner engine matched to the linearizations' character (monotone and
clean: \code{bsHe94b}; degenerate: the IPM adapter).

\subsection{\code{semismoothNewtonSolve} --- direct mixed-NCP Newton}
\label{sec:ssn}

The direct solver for the nonlinear mixed NCP, the library's core case.
Instead of projecting onto $K$, it compiles $K$ into a nonsmooth equation:
with an NCP function $\varphi$ ($\varphi(a,b) = 0 \iff a \ge 0,\ b \ge 0,\
ab = 0$),
\[
  \Phi(z) \;=\; \bigl[\, H(x,y)\;;\ \varphi(y_i, G_i(x,y)) \,\bigr] \;=\; 0
\]
is solved by generalized-Jacobian (semismooth) Newton, globalized by a
directional Armijo search on $\Psi = \tfrac12\|\Phi\|^2$ with a three-tier
direction ladder: the Newton direction; on a singular factorization or a
failed descent test ($\nabla\Psi \cdot d \le -\rho\|d\|^{p}$), a
Levenberg--Marquardt direction; as a last resort, the negative gradient.
There is no Josephy--Newton nesting and no smoothing parameter: unlike
\code{smoothingnewton} (Section~\ref{sec:blocks}), which needs a
$\mu$-continuation precisely because plain Newton stalls at the exact
Fischer--Burmeister kink, the semismooth method operates at the kink
itself.

\begin{lstlisting}[language=C++]
VIResult semismoothNewtonSolve(const VIModel& model, const VectorXd& z0,
                               const SemismoothNewtonParams& params = {},
                               const IterationLogger& logger = {});
\end{lstlisting}

Its two extension points:
\begin{itemize}
\item \code{NcpFunctionPair} --- the NCP function and its
  C-subdifferential diagonals, as a pair. Ready-made:
  \code{penalizedFischerBurmeisterPair(lambda = 0.8)} (the default;
  penalization repairs plain FB's documented weaknesses) and
  \code{fischerBurmeisterPair()} (the classical/restart option).
  \code{lambda} must lie in $(0,1)$ and should never go below $0.5$.
\item \code{JacobianFn} --- an analytic $F'$ hook; empty means the 4th-order
  FD Jacobian. For affine problems bind the constant matrix and the FD cost
  vanishes.
\end{itemize}

\code{SemismoothNewtonParams} adds the descent-test constants (\code{rho},
\code{pExp} $> 2$), the directional Armijo block, the nonmonotone
line-search memory (\code{nonmonotoneMemory}; 1 = monotone, the citable
theory; 4 = the production value that escapes merit plateaus), and the LM
damping scale.

Contract points that matter in practice: iterates are \emph{not} confined to
$K$, so $H$ and $G$ must tolerate evaluation at arbitrary points --- a trial
point where \code{evaluateF} throws counts as merit $+\infty$ \emph{during
the line search only} (the step backtracks away from domain violations); at
accepted iterates a non-finite value still throws. Termination and
\code{VIResult::residual} use the library-standard squared natural residual
(comparable across engines, and immune to a documented small-$\lambda$
merit-stop pitfall), while the line search steers by $\Psi$. The returned
point is the best visited, not the last.

\emph{Choose it} for genuinely nonlinear mixed NCPs --- it is the engine
built for exactly that shape --- and expect reported stalls rather than
convergence on nonmonotone problems far from a solution; the composition
below addresses that case.

\subsection{\code{alternatingChainSolve} --- the nonmonotone composition}
\label{sec:altchain}

For nonmonotone mixed NCPs where \emph{no} single engine converges, the
library composes two into rounds of
\[
  \text{project onto } K
  \;\longrightarrow\; \text{globalize}
  \;\longrightarrow\; \text{finish},
\]
with best-point memory across all stages.

\begin{lstlisting}[language=C++]
using StageSolver = function<VIResult(const VectorXd& start)>;
VIResult alternatingChainSolve(const VIModel& model, const VectorXd& z0,
                               const StageSolver& globalizer,
                               const StageSolver& finisher,
                               const AlternatingChainParams& params = {},
                               const ChainStageLogger& logger = {},
                               const Projector& projector = {});
\end{lstlisting}

The caller binds complete engines into the two \code{StageSolver} extension
points --- in the validated configuration, \code{solveVI} with the IPM
inner solver (under the stall cutoff) as globalizer and
\code{semismoothNewtonSolve} (nonmonotone memory 4) as finisher. The chain
itself: recomputes the squared natural residual for every stage result with
its \emph{own} fixed functional (stage solvers' conventions cannot skew the
comparison); keeps and finally returns the best point visited; treats a
stage that \emph{throws} as a stalled stage, not a failure (the round
continues with its other stage from the last good point --- the same
treatment PATH applies to subsolver failure); and, on a round that fails to
improve the best residual (\code{improveFactor}, default 1.0 = any strict
improvement continues, \code{roundsMax} alone bounds cost), restarts the
next round from the best point perturbed componentwise by
$U[-1,1] \cdot \code{perturbScale} \cdot \sqrt{\code{bestMag}}$ --- a
perturbation proportional to the current error, hence vanishing as the
chain approaches a solution, drawn from a fixed-seed generator so runs stay
reproducible. \code{perturbScale = 0} (default) disables perturbation, and
stagnation then ends the chain. The rationale for each element, with the
run evidence that forced it, is the subject of Section~\ref{ch:testing}'s
deploy case study and Chapter~\ref{ch:math}.

\emph{Choose it} when the problem is (or may be) nonmonotone and single
engines stall: it is bounded, reproducible, and exact in its convergence
reporting, and it is the only configuration in this library that has solved
the deployment game.

\subsection{Building blocks}
\label{sec:blocks}

Reusable pieces beneath the engines, each independently tested:

\begin{itemize}
\item \code{armijo.hpp} --- two backtracking searches on a caller-supplied
  scalar merit: \code{armijoLineSearch} (value-form test
  $\varphi(\alpha) \le (1 - c\,\alpha)\varphi_0$; no gradient needed) and
  \code{armijoLineSearchDirectional} (the classical
  $\varphi(\alpha) \le \varphi_0 + c\,\alpha\,\varphi'_0$ form the
  semismooth theory is stated for; pass a nonmonotone baseline for a
  nonmonotone search). On failure both return the best $\alpha$ tried with
  \code{accepted = false} --- the caller decides.
\item \code{levenbergmarquardt.hpp} --- the damping primitives
  \code{levenbergMarquardtDamp} ($J^\top J + \lambda I$; rectangular $J$
  accepted) and \code{levenbergMarquardtUpdate} (the $\lambda$ policy), plus
  \code{levenbergMarquardtSolve}, a self-contained nonlinear least-squares
  solver on those primitives and the FD Jacobian.
\item \code{fdjacobian.hpp} --- \code{centralDifferenceJacobian}: 4th-order
  central differences, per-component steps snapped to exactly representable
  widths; the real-arithmetic analogue of the Octave complex-step Jacobian.
\item \code{dampednewton.hpp} --- \code{dampedNewtonSolve}: a
  Gauss--Newton/gradient blend for least squares (the \code{dNewton.m}
  port), used by the smoothing solvers.
\item \code{smoothingnewton.hpp} --- the smoothed-FB reformulation solved by
  damped Newton, in two variants (merit-driven $u$, and
  \code{smoothingContinuationSolve} with $\mu$ as an outer homotopy
  parameter). Kept as the smoothing-family reference point; for mixed NCPs
  the semismooth solver (Section~\ref{sec:ssn}) supersedes it in practice
  because it needs no continuation schedule.
\end{itemize}

\section{Choosing an engine}
\label{sec:choosing}

Table~\ref{tab:engine-choice} reduces the catalog to the questions that
decide the choice, and the library also mechanizes the choice
(\code{chooseengine.hpp}). Three layers, separately usable:
\code{probeMonotone(M)} decides ``is $\mathrm{sym}(M)$ PSD?'' with a single
shifted Cholesky attempt rather than an eigensolve; the \emph{pure} function
\code{chooseEngine(ProblemTraits)} encodes the table over
$\{$dimension, \code{numFree}, affine?, orthant $K$?, monotone?, warm
start?$\}$, deterministic and unit-tested rule by rule; and two executors
run the choice --- \code{solveAffineAuto} (probe, pick
IPM/\code{bsHe94b}/chain, run; structured Newton factories pass through to
the IPM path) and \code{solveModelAuto} (semismooth first; on a reported
stall or a runtime guard, fall back to the alternating chain --- the
empirical policy of Chapter~\ref{ch:testing}). Every choice and fallback is
announced through a \code{ChoiceLogger} hook, never taken silently. Two
deliberate design points: there is no degeneracy proxy, because the IPM is
already the choice for every monotone cold start and
degeneracy-insensitivity is exactly why; and warm-start status is a
\emph{declared} flag, since only the caller can know whether its start is a
previous solution.

\begin{table}[htbp]
\centering
\small
\begin{tabular}{p{0.42\textwidth} p{0.50\textwidth}}
\toprule
\textbf{Situation} & \textbf{Engine} \\
\midrule
Affine, monotone, moderate conditioning; or warm start available; or $K$ is
not the mixed orthant (e.g.\ ellipsoid)
  & \code{bsHe94b} \\
\addlinespace
Affine, monotonicity doubtful, bad start, or matrix-free required
  & \code{solodovSvaiter}, usually as phase 1 of \code{chainedSolodovHe} \\
\addlinespace
Affine, monotone, but degenerate / near-tied / large (KKT of a convex QP)
  & \code{mehrotraIpm}; add a structured \code{NewtonSolverFactory} when the
    problem has one \\
\addlinespace
Nonlinear, smooth, $K$ arbitrary
  & \code{solveVI} + inner engine matched to the linearizations \\
\addlinespace
Nonlinear mixed NCP (the core case)
  & \code{semismoothNewtonSolve}; analytic \code{JacobianFn} if available \\
\addlinespace
Nonmonotone; single engines stall
  & \code{alternatingChainSolve} (IPM globalizer + semismooth finisher;
    enable \code{perturbScale} for stagnation) \\
\bottomrule
\end{tabular}
\caption{Engine selection at a glance.}
\label{tab:engine-choice}
\end{table}

\section{The network layer}
\label{sec:network}

The flow-planning application is contained in its own static library,
\code{vincpnet} (\code{network/include}, \code{network/lib}, namespace
\code{VINCP::Network}), layered strictly on top of \code{vincp}: the core
library knows nothing about networks. It is both a production solver for the
distribution-planning QP and the library's largest worked example of the
``data plus builders'' pattern. The mathematical formulation is
Section~\ref{sec:a-flowqp}; here is the API.

\subsection{Instances and plans}

\code{Instance} (\code{instance.hpp}) is the problem datum: per-node supply
caps $C_i$, demands $D_i$, priorities $P_i$; the full cost matrix $c_{ij}$
(positive diagonal --- self-supply is routed and priced, not free); the
ton-mile budget \code{tonMileLimit} $L$; and optional display geometry
(\code{xCoord}/\code{yCoord}, \code{labels}), validated only when present so
hand-built abstract instances may omit them. \code{makeRandomInstance(profile,
seed)} generates instances from a validated \code{InstanceProfile}: node-class
counts (supply-only / both / demand-only / transit), tonnage and priority
ranges, and the geometry family --- \code{laydownType 0} (random square,
near-metric costs) or \code{1} (banded rectangles, near-tied costs; the
deliberately degenerate family).

\code{Plan} (\code{plan.hpp}) is the answer shape shared by every planner:
supplied tonnage $S_i$, resupply $R_i$, and the flow matrix $f_{ij}$;
\code{checkPlan} returns a \code{PlanViolations} breakdown (balance,
delivery, capacity, negativity, idle resupply, budget) and
\code{shortfallObjective} evaluates the dimensionless objective
$\theta = \sum_i P_i \bigl((\mathit{target}_i - R_i)/D_i\bigr)^2$. One
recorded caution: $\theta$ is a weighted sum of squared \emph{fractions} ---
$O(1)$ to a few hundred --- not a tonnage; read delivered/unmet tons
separately.

\subsection{Baseline planners}

Three non-optimizing planners provide baselines, starts, and GUI overlays:
\code{greedyPlan} (\code{greedy.hpp}) --- the rationing-aware greedy
heuristic, which also \emph{calibrates the budget}: a newly constructed
\code{Instance} has \code{tonMileLimit = 0} and the solvers throw on a
non-positive budget, so the standard sequence sets \code{inst.tonMileLimit =
greedyResult.suggestedLimit} ($\approx 80\%$ of greedy usage) first;
\code{gravityPlan} (\code{gravity.hpp}) --- the dense cost-blind
proportional baseline; and the 2-exchange engine (\code{swap.hpp}) ---
\code{bestSwapAtNode} / \code{bestSwap} / \code{swapNodeToLocalOptimum} /
\code{swapToLocalOptimum} over arcs that already carry flow (supplied and
resupply are provably invariant under a swap; only routing changes).

\subsection{Reduction, screening, and the LCP}

\code{reduction.hpp} turns the node problem into a source--sink pair
problem: \code{computeShortestRoutes} (all-pairs cheapest routes, including
the self-supply detour \code{selfDistance}/\code{selfVia}), then
\code{makeReducedProblem} with a \code{ScreenParams} --- a count rule
(\code{maxSourcesPerSink}: keep the $k$ cheapest sources per sink) and a gap
rule (\code{gapFraction}: keep sources within $(1+\mathit{gap})$ of the
cheapest), applied as a union; both zero keeps every pair.

\code{buildFlowLcp} (\code{flowlcp.hpp}) assembles the monotone KKT-LCP over
the kept pairs, packed \emph{sink-major}: $z = [\,t\ (\text{pair flows, each
sink's pairs contiguous}) \mid \mu\ (\text{per-source capacity duals}) \mid
\lambda\ (\text{budget dual})\,]$, all $\ge 0$ (\code{numFree = 0});
\code{unpackFlowLcp} routes the pair solution back along the shortest routes
into a node-level \code{Plan}. The tie-break term (R4) enters here:
\code{epsilon} $< 0$ selects \code{defaultTieBreakEpsilon(inst)}.

\code{makeFlowNewtonFactory(lcp)} (\code{flownewton.hpp}) is the structured
\code{NewtonSolverFactory} of Section~\ref{sec:ipm}: per-sink
Sherman--Morrison on the rank-1 $t$-blocks plus a small SPD Schur complement
on the $\sim(\text{numSources}+1)$ duals, replacing the dense $O(d^3)$
factorization with $O(\text{pairs})$-scale work per IPM iteration. Its
algebra was verified symbolically with Maxima \cite{maxima}
(\code{network/doc/ns2-newton-check.mac}).

For tiny instances only, \code{oracle.hpp} builds and solves the full
\emph{unreduced} formulation (\code{solveOracle}, capped at
\code{maxNodes = 8}) --- the ground-truth cross-check for the reduction.

\subsection{\code{solveFlowPlan}: the production entry point}

\begin{lstlisting}[language=C++]
FlowPlanResult solveFlowPlan(const Instance&, const FlowPlanParams& = {});
\end{lstlisting}

\code{FlowPlanParams} bundles the pipeline's controls: the engine selection
(\code{engine} $\in$ \{\code{bshe94b}, \code{chain}, \code{ipm},
\code{ssn}\}, plus \code{ipmNewton} $\in$ \{\code{dense}, \code{flow}\} and
\code{newtonCheckTol} for the IPM), the chain's phase-1 controls
(\code{roughMagTol}, \code{roughIterMax}), tolerances (\code{magTol},
\code{iterMax}, \code{iterFreq}), the screen (\code{maxSourcesPerSink},
\code{gapFraction}), the exactness-restoring certificate loop
(\code{maxCertificateRounds}, \code{certificateSlack}), and the tie-break
\code{epsilon}. Remember the engine-dependent meaning of \code{iterMax}:
under \code{ipm}/\code{ssn} it counts factorizations (set low hundreds),
under the projection engines iterations (hundreds of thousands).

\code{FlowPlanResult} returns the \code{Plan} in real tons, the objective
\code{shortfall} ($\theta$), \code{tonMilesUsed}, the budget shadow price
$\lambda$ (real units), the raw \code{vi} result (\emph{scaled} units,
squared residual), the kept/total pair counts, \code{certificateRounds},
and \code{certifiedP}. Two reading rules: trust the plan as optimal only
when \code{certifiedP} is true (and \code{vi.converged}); and remember the
unit split --- tolerances and \code{vi.residual} are in the solver's
nondimensionalized units (tons scaled by the largest demand, miles by the
largest cost: an exact unit change the projection engines need to converge
at a sane rate), while the plan and shadow price come back in real units.

Screening and exactness compose as follows: the screen only pre-trims
pairs; each certificate round checks the excluded pairs' reduced costs and
pulls violators back in, so the final answer is exact whenever the loop
certifies. Keep-all with \code{maxCertificateRounds = 0} is exact by
construction --- nothing excluded, nothing to certify --- and the
\code{ipm} + \code{ipmNewton = flow} combination exists precisely to make
that affordable at scale. The production defaults by geometry, from the
performance study: banded or large $\Rightarrow$ \code{ipm} + \code{flow} +
no screen; small spread-geometry instances keep screened \code{bshe94b}.

\subsection{Configuration files and the benchmark harness}

\code{config.hpp} parses \code{key = value} text (\code{\#} comments;
duplicate keys, empty keys or values, and --- via
\code{requireAllConsumed} --- unknown or misspelled keys are hard errors, so
a mistyped key can never silently have no effect). \code{applyProfileConfig}
consumes the \code{profile.*} keys (all \code{InstanceProfile} fields);
\code{applyFlowPlanConfig} the \code{solver.*} and \code{screen.*} keys (all
\code{FlowPlanParams} fields). The benchmark driver
(\code{network\_benchmark <path-to-cfg>}, a plain Release executable) adds
\code{benchmark.name}, \code{benchmark.instances} (default 5), and
\code{benchmark.seedBase}. Two shipped configurations serve as templates:
\code{benchmark-example.cfg} (screened banded run) and
\code{ns3-keepall.cfg} (the keep-all \code{ipm}+\code{flow} production
default).

\subsection{The instance viewer}

\code{network\_viewer} (\code{network/gui}, Qt6, behind
\code{option(VINCPNET\_BUILD\_GUI)} with a self-disabling subdirectory when
Qt is absent --- the solver build never depends on Qt) is the interactive
companion: instance generation controls, the map canvas
(\code{FlowPlanView}: nodes at true coordinates, pan/zoom, closest-links and
plan overlays, node popups, swap gestures), a cost histogram, node rankings,
and a status line showing $\theta$, ton-miles, and delivered/unmet tons for
whichever plan is displayed.

\section{Bringing a new problem}
\label{sec:new-problem}

A new problem enters the library as \emph{data plus a \code{VIModel}
builder} --- never as copied solver or harness code. The pattern, proven
twice by the GAMS translations, is a recipe worth stating step by step,
because fidelity depends on its details.

\subsection{The GAMS translation recipe}
\label{sec:gams-recipe}

Given a GAMS MCP --- a \code{Model} statement pairing equations with
variables --- the C++ \code{VIModel} is constructed so it is \emph{the same
MCP}, mechanically:

\begin{enumerate}
\item \textbf{Pack the nonnegative block in the Model statement's pairing
  order.} Every GAMS \code{Positive Variable} becomes a slice of $y$;
  compute the offsets once as named \code{constexpr} indexers
  (\code{prAt(r)}, \code{flowRAt(r,m,k)}, \ldots) and use them everywhere
  --- they are also what makes diagnostic output legible.
\item \textbf{Pair \code{=g=} rows with their variables as $G$.} Each GAMS
  inequality \code{lhs =g= 0} paired with variable $v$ becomes
  $G_{[v]} = \mathit{lhs}$, transcribed literally --- including every
  $\varepsilon$-guard the source carries.
\item \textbf{Pair \code{=e=} rows with free multipliers as $H$.} Equality
  constraints and their (sign-unrestricted) multipliers form the free block
  $x$; intermediate GAMS variables defined by equalities may also be placed
  here when they are provably interior (document the argument --- e.g.\ a
  variable defined as $\varepsilon$ plus a sum of nonnegatives is strictly
  positive, so its GAMS \code{Positive} bound never binds and the free
  treatment is exact).
\item \textbf{Drop only what is provably redundant, and say so.} A box bound
  implied by another constraint (e.g.\ an individual bound implied by a
  budget row) may be dropped, with the multiplier-absorption argument
  recorded in a comment.
\item \textbf{Reproduce the \code{.L} levels as the start.} GAMS initial
  levels are part of the model's intent --- for the deployment game,
  interior starts are \emph{essential} (the survivor function has zero slope
  at zero) --- so \code{initialPoint()} transcribes them, including their
  deliberate infeasibilities.
\item \textbf{Print the GAMS output parameters.} A report-only check that
  reproduces the \code{.gms} file's own \code{Display} quantities (option
  probabilities, payoffs, \ldots) lets a converged run be compared against
  the GAMS listing by eye before any automated check is fixed to it.
\item \textbf{Check against verified answers.} Once the author confirms a
  solution matches the reference system, fix the test to that answer; for
  games with multiple equilibria, check against a \emph{roster} of verified
  equilibria and have the check print any converged non-roster point as a
  paste-ready roster candidate (Chapter~\ref{ch:testing}).
\end{enumerate}

Cross-validation closes the loop: two independent engine classes converging
to the same author-verified point is strong evidence the transcription is
faithful --- a sign error almost anywhere would split them.

\subsection{The GMS front end: the recipe, mechanized}
\label{sec:gms-frontend}

The \code{gms/} layer (library \code{vincpgms}, namespace
\code{VINCP::Gms}) performs the translation recipe mechanically for the
subset of the GMS language\footnote{``GAMS'' is a registered trademark of
GAMS Development Corporation. This front end is not endorsed or certified
by GAMS Development Corporation, and the subset it accepts is
incompatible with most of the GAMS modeling language; every source file
in the layer carries the full disclaimer.} that the project's model
corpus uses. The subset is defined by a construct-by-construct census of
the six corpus files (\code{doc/2026-07-08-gams-subset-census.md}); the
front end accepts exactly that footprint and throws
\code{std::invalid\_argument}, with a file, line, and column, on anything
outside it --- a sixth file's novelty fails loudly rather than parsing
wrongly.

Four layers, each with its own test suite:

\begin{itemize}
\item \textbf{Lexer and parser} (\code{gmslexer}, \code{gmsparser}):
  tokens carry the original spelling plus a lower-cased key (GMS
  identifiers are case-insensitive, and the corpus depends on this);
  \code{\$include} splices, \code{\$macro} expands at token level with
  rescanning for nested macros; comments count only in column one. The
  parse tests assert the census counts on every corpus file, and a
  canonical echo (\code{gmsecho}) gates the round trip on idempotence.
\item \textbf{Database and evaluator} (\code{gmsdatabase},
  \code{gmseval}): \code{buildGmsDatabase} walks the program in statement
  order --- declarations create symbols, tables and data blocks fill
  dense arrays, assignments evaluate immediately with indexed left-hand
  sides looping over their domains, equation definitions are stored and
  checked symbolically. \code{Solve} is recorded, not executed, so
  assignments after it evaluate at the initial levels; the tests pin
  derived parameters to hand-computed values.
\item \textbf{Model builder} (\code{gmsmcp}): \code{buildGmsMcp} turns
  one \code{Model} statement into a \code{VIModel} by the pairing rules
  of Chapter~\ref{ch:math} (the variable's kind decides its block; the
  relation is checked for consistency only), packs $z_0$ from the
  \code{.L} levels, and \emph{surfaces} \code{.UP} bounds without
  applying them. The returned closures read the database live, so the
  database must outlive the model.
\item \textbf{Write-back} (\code{applyMcpSolution},
  \code{rerunPostSolveAssignments}): a solution vector lands back in the
  variables' \code{.L} arrays, and the post-\code{Solve} assignments are
  re-executed so the file's own report parameters are computed at the
  solved point, comparable line by line against a GAMS listing.
\end{itemize}

The canonical usage is six calls:

\begin{lstlisting}[language=C++]
Program prog   = parseGmsFile("model.gms");        // throws outside the subset
GmsDatabase db = buildGmsDatabase(prog);           // data + validated equations
GmsMcp mcp     = buildGmsMcp(db, "modelname");     // VIModel + z0 + slots
VIResult r     = solveModelAuto(mcp.model, mcp.z0);
applyMcpSolution(db, mcp, r.z);                    // z -> the .L levels
rerunPostSolveAssignments(db, prog);               // report parameters at z
\end{lstlisting}

The front end's role is reference, prototyping, and cross-check: four of
the five corpus models reproduce their GAMS solutions through it
(Chapter~\ref{ch:testing}, Section~\ref{sec:t-gms}), and a hand-built
production twin can be checked row by row against the parsed model. It
is not the production path: interpreted evaluation carries measured
costs (Section~\ref{sec:t-gms-cost}) that hand-built structures avoid,
and the decision of record keeps production problems on the recipe
above.

\subsection{The multi-engine test harness}

\code{test/mcpengines.hpp} binds one \code{VIModel} to named solver rows so
an acceptance test is data plus assertions: \code{makeMcpEngineRows(model,
engines, params)} produces rows from \code{\{Ssn, JnHan, JnHe, JnIpm\}},
\code{makeSsnRow} adds parameter-variant rows,
\code{countConvergedRows(rows, z0, extraChecks)} runs each through the
shared \code{runCase} driver (per-row timing and statistics, throws absorbed
and reported, output flushed per row). Because nonmonotone problems make
individual engines fail legitimately, the checks are expressed over the row
results (``at least $k$ rows pass''), not per row.

\subsection{The network walkthrough}

For the flow-planning layer the same philosophy takes pipeline form; the
canonical usage is five calls:

\begin{lstlisting}[language=C++]
Instance inst = makeRandomInstance(profile, seed);   // 1. generate (or hand-build)
GreedyResult g = greedyPlan(inst);                   // 2. baseline + budget
inst.tonMileLimit = g.suggestedLimit;                //    calibrate (mandatory)
FlowPlanParams p;                                    // 3. choose engine/screen
p.engine = "ipm"; p.ipmNewton = "flow";              //    (banded/large default)
FlowPlanResult r = solveFlowPlan(inst, p);           // 4. reduce+solve+certify
// 5. read r.plan (real tons) -- but only trust it when
//    r.certifiedP && r.vi.converged.
\end{lstlisting}

Steps 2--3 are the ones most often done incorrectly: the budget calibration
is not optional (an uncalibrated instance throws), and the engine choice
changes the meaning of \code{iterMax}. Everything between reduction and
unpacking --- screens, the LCP build, the certificate loop,
nondimensionalization --- is internal to \code{solveFlowPlan}; the
individual stages remain public (Section~\ref{sec:network}) for tests and
custom pipelines.

\section{Running and writing tests}
\label{sec:tests}

Section~\ref{sec:getting-started} gave the one-line commands; this section
is the working reference: how to run subsets, how to read a suspicious
run, and the rules a new test must follow. The taxonomy of what the suite
checks (unit / protocol / acceptance) and the evidence it produced are in
Chapter~\ref{ch:testing}; this section is about mechanics.

\subsection{Running tests}

\code{ctest} filters by regular expression over test names, and the names
are chosen to make the filters useful: every network-layer suite name
begins with \code{Network}, so \code{ctest -R "\^{}Network"} runs exactly
that layer, \code{ctest -R "\^{}NetworkFleet"} the fleet subset, and
\code{ctest -R MehrotraIpm} one engine's suite. Always pass
\code{--output-on-failure}. In CLion the aggregate targets
(\code{run\_all\_tests}, \code{run\_engine\_tests},
\code{run\_network\_tests}) execute ctest during the \emph{build} step:
use Build (the hammer) on them, never Run --- a custom target has no
executable, and Run reports that none is specified. The native ``All
CTest'' run configuration is the alternative when the Run button is
wanted.

Two mechanics deserve emphasis because each cost a debugging session.
First, adding a test file or target requires a CMake \emph{reload} before
the build, not merely a recompile; editing only existing sources needs a
plain rebuild. Second, a stale build announces itself in the ctest
header: the suite prints its total test count, and a run whose count does
not match the source tree did not build what is in the editor. A
2026-07-08 run reported a failure that no current source could produce;
the count (177 where the tree held 182) identified the binaries as stale
before any code was inspected.

The benchmark programs (\code{network\_benchmark},
\code{fleet\_benchmark}) are deliberately \emph{not} tests: they are
plain executables, run manually, and their timings are meaningful only in
a Release configuration. Both print flushed per-iteration and per-round
progress, so a run killed mid-solve still reports the dimension and
per-iteration cost it was attempting --- the property that made the
bounded probes of Chapter~\ref{ch:testing} cheap to run.

\subsection{Writing tests}

A new test is a \code{.cpp} in \code{test/} (core) or
\code{network/test/} (network layer), registered by one line ---
\code{vincp\_add\_gtest(name)} or \code{vincpnet\_add\_gtest(name)} ---
which creates the executable, links it, and enrolls it in discovery and
the aggregate targets. Network-layer suite names must begin with
\code{Network}, or the aggregate regex will silently exclude them. All
dimensions and tolerances are named constants; solver cases run through
the \code{runCase}/\code{CheckFn} harness of
Section~\ref{sec:core-types} so timing, statistics, and throw-reporting
are uniform.

Beyond the mechanics, four design rules from this project's experience:

\begin{enumerate}
\item \textbf{Choose the comparison bar by what it can prove.} Three bars
  recur, in increasing strictness. A \emph{backward-error} bar
  ($\|Kd - \mathit{rhs}\| / (\|K\|_\infty \|d\| + \|\mathit{rhs}\|)$)
  is valid at any conditioning and is the correct parity check for a
  reimplemented linear solve. A \emph{direct} comparison of solutions or
  steps carries the problem's conditioning and is asserted only on
  well-scaled systems --- the parity fixtures are nondimensionalized
  exactly as the production path is, for this reason. An \emph{exact}
  (bit-for-bit) equality gate is reserved for refactorings whose claim is
  ``the old path is unchanged'': the factory extension point and the
  matrix-free overload of \code{mehrotraIpm} are each guarded by a test
  asserting exact equality of iterates, iteration count, and residual
  against the historical path.
\item \textbf{Test protocols with deliberately failing fakes, not with
  hard problems.} Rescue paths, drift guards, stall cutoffs, and chain
  control flow are driven by injected components through the same
  extension points production code uses. Chapter~\ref{ch:testing}
  documents why the alternative fails, including a guard that a
  plausible-looking natural problem had never actually triggered.
\item \textbf{A randomized test must assert draw-independent
  invariants.} The convention \code{makeSeed(0)} draws a fresh seed from
  the clock each run (and prints it, so a failure is reproducible); such
  a test is a per-run fuzzer, and its assertions must hold on
  \emph{every} draw. The recorded counterexample (2026-07-08): a
  cross-engine check compared two converged solution \emph{points} at
  $10^{-5}$, and an unlucky draw produced a QP whose solution set was
  nearly flat --- both engines were correct, and the points legitimately
  differed by $1.4 \times 10^{-5}$. The assertion now compares the
  objective value, which is constant across a convex QP's solution set,
  with only a gross-disagreement bar on the points. Fixed-seed tests
  assert exact expectations; reroll tests assert invariants.
\item \textbf{Verify intricate algebra symbolically before implementing
  it.} The structured Newton factories were derived on paper, checked in
  Maxima \cite{maxima} (exact rational and symbolic arithmetic, one
  \code{PASS} line per claim, a completeness counter so a silently
  skipped check fails the script), and only then transcribed; the C++
  parity tests then verify the transcription, not the mathematics. The
  scripts ship in \code{network/doc/} and the corresponding sections of
  Chapter~\ref{ch:math} cite them check by check.
\end{enumerate}

Finally, the license banner convention applies to test sources like any
other file, and a check function's report string should state what was
compared and against what bar --- the failure message is the first thing
read six months later, usually by someone who has forgotten the test
exists.

\section{Extending the library}
\label{sec:extending}

\paragraph{Adding an engine.} Implement the calling convention of its group
(the affine signature, or \code{VIModel} in / \code{VIResult} out); validate
inputs up front (\code{std::invalid\_argument}); observe the
squared-residual and stall-vs-throw contracts; report through the shared
logger types. Provide: a params struct with documented defaults, a header
contract comment (they are the API of record --- this chapter is derived
from them), unit tests on constructed-solution problems, and --- if the
engine participates in drivers --- a \code{make*Solver} adapter and a
comparison row in the cross-engine tests. Cite sources in the header
(every engine in the library carries its literature).

\paragraph{Adding an NCP function.} Provide an \code{NcpFunctionPair}: the
elementwise value and the C-subdifferential diagonals $(d_a, d_b)$,
including a kink rule consuming the limiting-direction argument. Note the
subtlety recorded from experience: the kink indicator is defined in the full
$z$-space (length $n+m$), so pure-LCP tests ($n = 0$) cannot detect a
$y$-space/$z$-space confusion --- include a mixed test that hits an exact
kink.

\paragraph{Adding a Newton factory.} Honor the \code{NewtonSolverFactory}
contract exactly ($K = M + \mathrm{blockdiag}(\mathit{reg}\, I_n,
\mathrm{diag}(\mathit{sOverY}))$, rebuilt per call, rescue protocol carried
through \code{freeRegularization}); develop with \code{newtonCheckTol > 0};
test with backward-error parity bars against the dense default
($\|Kd - \mathit{rhs}\|$ relative bars always; direct step parity only on
well-scaled systems, since forward error carries $\mathrm{cond}(K)$). When
the algebra is intricate, verify it symbolically with Maxima
\cite{maxima} first, as the flow factory does: its shipped script checks
every formula before any C++ ran.

\paragraph{The JNI boundary (planned).} Eigen is column-major by default; a
row-major buffer arriving from Java must be mapped with an explicit
\code{RowMajor} map or it transposes silently. This is the one recorded
hazard for the future Java binding.
% ----------------------------------------------
% Copyright Ben Paul Wise. All Rights Reserved.
% ----------------------------------------------
